% !TEX root = ../main.tex
% Flush the baseline figures of Section 6 before this section starts.
\FloatBarrier
\section{Quantizing the agent-to-server transmissions}
\label{sec:quantization}

FedNPG-ADMM reduces communication by shrinking \emph{what} is transmitted from
a matrix to two vectors. A second axis, which the original work does not
explore, is to shrink \emph{how precisely} those vectors are transmitted. The
original paper reports overhead at 64 bits per parameter, and our byte counter follows
it, although the vectors our implementation exchanges are \texttt{float32}.
Compression ratios below are therefore given against both references: the
paper's 64-bit accounting and the 32 bits actually sent. If performance
survives at 8 or 4 bits, the saving compounds with the $\Oh(d^2) \to \Oh(d)$
one.

This is the novel contribution of the project. The precedent is Q-GADMM
\citep{elgabli2021qgadmm}, which \citet{lan2023fednpgadmm} cite themselves, and
which quantizes the vectors an ADMM scheme transmits in the \emph{supervised}
setting. We carry that idea into the reinforcement learning setting, where the
transmitted vectors are gradient and consensus directions computed from a
non-stationary sampling distribution. We found no prior work quantizing
agent-to-server transmissions for federated policy gradient with
continuous-control experiments. The nearest results are for compressed temporal
difference learning with linear function approximation
\citep{mitra2024compressedtd}, which is theory-only and not policy gradient.

\subsection{Design}

A quantizer acts on each vector an agent sends to the server. It maps the
vector to $b$ bits per entry, reconstructs it on arrival, and declares the bit
width the transmission would have needed. It sits entirely on the
communication channel, so the algorithm itself is unchanged.

Two properties make the experiment measure something. The degraded tensor is
what the server aggregates, because downstream code uses the return value of
\texttt{send} and never the original, so the quantization error enters the
update to $\bthe$. And the declared bit width is what the
byte meter uses, so the cost axis moves with the fidelity. Per iteration, with
$R$ ADMM rounds, $N$ agents and $b$ bits, the uplink is
\begin{equation}
  \text{bytes} \;=\; N \,(1 + R)\, d \, b / 8,
  \label{eq:quantbytes}
\end{equation}
since $\bg_i$ is sent once per iteration and $\by_i$ once per round. Measured
byte counts match \eqref{eq:quantbytes} exactly; Figure~\ref{fig:quantoverhead}
shows what they add up to over a full training run.

\paragraph{Stochastic uniform quantization.} For a tensor $x$ at $b$ bits, with
$L = 2^b - 1$ levels and step $s = (\max x - \min x)/L$, the value is placed at
$\lfloor u \rfloor + \mathrm{Bernoulli}(u - \lfloor u \rfloor)$ where
$u = (x - \min x)/s$. Rounding stochastically instead of to nearest makes the
estimator unbiased, $\Exp[\hat{x}] = x$, with per-element variance at most
$s^2/4$. That property is the reason to prefer it here: a deterministic
rounding error is systematic and would accumulate across iterations, whereas an
unbiased one leaves the server receiving the correct gradient and consensus
direction in expectation. The quantizer draws from a dedicated
\texttt{torch.Generator} rather than the global stream, so runs at different
bit widths see identical action-sampling sequences and differ only in the
quantization.

\subsection{Results}

The main experiment is a five-seed grid: stochastic quantization at 8 and 6
bits, for $N \in \{2,4,8\}$, on all three environments (90 runs). One agent has
no uplink to quantize, so $N=1$ is unquantized throughout. Each environment
runs in its own configuration (Section~\ref{sec:implementation}), and the
quantized runs use the same settings and seeds as the unquantized runs they
are compared against: the single-agent Swimmer runs of the two code bases agree to the last
logged digit. Table~\ref{tab:quantgrid} and Figure~\ref{fig:quant} give the
results; a finer single-seed scan over bit widths at $N=8$ locates the
threshold more precisely.

\begin{table}[t]
  \centering
  \caption{Final reward under stochastic quantization, mean $\pm$ standard
  deviation across five seeds, computed as in Table~\ref{tab:reproduction}.
  The unquantized column exchanges \texttt{float32} values; 8 and 6 bits cut
  the uplink per iteration by $4\times$ and $5.3\times$ against it, or $8\times$
  and $10.7\times$ at the paper's 64-bit accounting. The unquantized column
  is the corresponding row of Table~\ref{tab:reproduction}.}
  \label{tab:quantgrid}
  \begin{tabular}{llrrr}
    \toprule
    Environment & $N$ & Unquantized & 8-bit & 6-bit \\
    \midrule
    \multirow{3}{*}{Swimmer-v4}
      & 2 & $126.0 \pm 1.3$ & $126.1 \pm 0.8$ & $128.9 \pm 7.4$ \\
      & 4 & $127.0 \pm 0.8$ & $127.1 \pm 0.5$ & $125.8 \pm 3.4$ \\
      & 8 & $128.4 \pm 0.8$ & $128.6 \pm 0.4$ & $127.8 \pm 0.7$ \\
    \midrule
    \multirow{3}{*}{Humanoid-v4}
      & 2 & $3391 \pm 796$ & $2956 \pm 593$ & $497 \pm 34$ \\
      & 4 & $5343 \pm 121$ & $5244 \pm 237$ & $676 \pm 73$ \\
      & 8 & $6329 \pm 334$ & $5997 \pm 259$ & $1514 \pm 526$ \\
    \midrule
    \multirow{3}{*}{Hopper-v4}
      & 2 & $1636 \pm 282$ & $2052 \pm 389$ & $1941 \pm 461$ \\
      & 4 & $2536 \pm 116$ & $2197 \pm 533$ & $2096 \pm 243$ \\
      & 8 & $2656 \pm 286$ & $2581 \pm 466$ & $2590 \pm 338$ \\
    \bottomrule
  \end{tabular}
\end{table}

\begin{figure}[tbp]
  \centering
  \pgfplotslegendfromname{quantlegend}\par\medskip
  %% Quantization: return over iterations, unquantized (float32), 8 and 6 bits,
%% one panel per environment (rows) and agent count N = 2, 4, 8 (columns).
%% Mean over five seeds; band spans min to max. Hopper uses its own
%% configuration, Swimmer and Humanoid the original one.
%% Data: figures/quant/quant_<env>_N<n>.dat (scripts/export_quant_curves.py).
%% The legend is stored as quantlegend; the figure places it with
%% \pgfplotslegendfromname{quantlegend}.
\begin{tikzpicture}
\begin{groupplot}[
    group style={group size=3 by 3, horizontal sep=0.85cm, vertical sep=0.75cm,
                 xlabels at=edge bottom, ylabels at=edge left,
                 yticklabels at=edge left},
    scale only axis, width=4.0cm, height=2.7cm,
    ymajorgrids=true, xmajorgrids=true, grid style=dashed,
    tick label style={font=\scriptsize},
    label style={font=\scriptsize},
    title style={font=\footnotesize, yshift=-0.5ex},
    ylabel style={align=center},
    xlabel={Iteration},
]

\nextgroupplot[xmin=0, xmax=1000,
    ymin=0, ymax=150,
    xtick distance=500,
    ytick distance=50,
    ylabel={Swimmer-v4\\return},
    title={$N=2$},
    legend to name=quantlegend, legend columns=3,
    legend style={font=\footnotesize, draw=none, /tikz/every even column/.append style={column sep=0.8em}}]
\addplot[name path=swimmerN2fp32lo, draw=none, forget plot] table[x=iteration, y=fp32_lo] {figures/quant/quant_swimmer_N2.dat};
\addplot[name path=swimmerN2fp32hi, draw=none, forget plot] table[x=iteration, y=fp32_hi] {figures/quant/quant_swimmer_N2.dat};
\addplot[colorA, opacity=0.15, forget plot] fill between[of=swimmerN2fp32lo and swimmerN2fp32hi];
\addplot[colorA, line width=0.8pt] table[x=iteration, y=fp32_mean] {figures/quant/quant_swimmer_N2.dat};
\addlegendentry{Unquantized (float32)}
\addplot[name path=swimmerN2b8lo, draw=none, forget plot] table[x=iteration, y=b8_lo] {figures/quant/quant_swimmer_N2.dat};
\addplot[name path=swimmerN2b8hi, draw=none, forget plot] table[x=iteration, y=b8_hi] {figures/quant/quant_swimmer_N2.dat};
\addplot[colorB, opacity=0.15, forget plot] fill between[of=swimmerN2b8lo and swimmerN2b8hi];
\addplot[colorB, line width=0.8pt] table[x=iteration, y=b8_mean] {figures/quant/quant_swimmer_N2.dat};
\addlegendentry{8 bits}
\addplot[name path=swimmerN2b6lo, draw=none, forget plot] table[x=iteration, y=b6_lo] {figures/quant/quant_swimmer_N2.dat};
\addplot[name path=swimmerN2b6hi, draw=none, forget plot] table[x=iteration, y=b6_hi] {figures/quant/quant_swimmer_N2.dat};
\addplot[colorC, opacity=0.15, forget plot] fill between[of=swimmerN2b6lo and swimmerN2b6hi];
\addplot[colorC, line width=0.8pt] table[x=iteration, y=b6_mean] {figures/quant/quant_swimmer_N2.dat};
\addlegendentry{6 bits}

\nextgroupplot[xmin=0, xmax=1000,
    ymin=0, ymax=150,
    xtick distance=500,
    ytick distance=50,
    title={$N=4$}]
\addplot[name path=swimmerN4fp32lo, draw=none, forget plot] table[x=iteration, y=fp32_lo] {figures/quant/quant_swimmer_N4.dat};
\addplot[name path=swimmerN4fp32hi, draw=none, forget plot] table[x=iteration, y=fp32_hi] {figures/quant/quant_swimmer_N4.dat};
\addplot[colorA, opacity=0.15, forget plot] fill between[of=swimmerN4fp32lo and swimmerN4fp32hi];
\addplot[colorA, line width=0.8pt] table[x=iteration, y=fp32_mean] {figures/quant/quant_swimmer_N4.dat};
\addplot[name path=swimmerN4b8lo, draw=none, forget plot] table[x=iteration, y=b8_lo] {figures/quant/quant_swimmer_N4.dat};
\addplot[name path=swimmerN4b8hi, draw=none, forget plot] table[x=iteration, y=b8_hi] {figures/quant/quant_swimmer_N4.dat};
\addplot[colorB, opacity=0.15, forget plot] fill between[of=swimmerN4b8lo and swimmerN4b8hi];
\addplot[colorB, line width=0.8pt] table[x=iteration, y=b8_mean] {figures/quant/quant_swimmer_N4.dat};
\addplot[name path=swimmerN4b6lo, draw=none, forget plot] table[x=iteration, y=b6_lo] {figures/quant/quant_swimmer_N4.dat};
\addplot[name path=swimmerN4b6hi, draw=none, forget plot] table[x=iteration, y=b6_hi] {figures/quant/quant_swimmer_N4.dat};
\addplot[colorC, opacity=0.15, forget plot] fill between[of=swimmerN4b6lo and swimmerN4b6hi];
\addplot[colorC, line width=0.8pt] table[x=iteration, y=b6_mean] {figures/quant/quant_swimmer_N4.dat};

\nextgroupplot[xmin=0, xmax=1000,
    ymin=0, ymax=150,
    xtick distance=500,
    ytick distance=50,
    title={$N=8$}]
\addplot[name path=swimmerN8fp32lo, draw=none, forget plot] table[x=iteration, y=fp32_lo] {figures/quant/quant_swimmer_N8.dat};
\addplot[name path=swimmerN8fp32hi, draw=none, forget plot] table[x=iteration, y=fp32_hi] {figures/quant/quant_swimmer_N8.dat};
\addplot[colorA, opacity=0.15, forget plot] fill between[of=swimmerN8fp32lo and swimmerN8fp32hi];
\addplot[colorA, line width=0.8pt] table[x=iteration, y=fp32_mean] {figures/quant/quant_swimmer_N8.dat};
\addplot[name path=swimmerN8b8lo, draw=none, forget plot] table[x=iteration, y=b8_lo] {figures/quant/quant_swimmer_N8.dat};
\addplot[name path=swimmerN8b8hi, draw=none, forget plot] table[x=iteration, y=b8_hi] {figures/quant/quant_swimmer_N8.dat};
\addplot[colorB, opacity=0.15, forget plot] fill between[of=swimmerN8b8lo and swimmerN8b8hi];
\addplot[colorB, line width=0.8pt] table[x=iteration, y=b8_mean] {figures/quant/quant_swimmer_N8.dat};
\addplot[name path=swimmerN8b6lo, draw=none, forget plot] table[x=iteration, y=b6_lo] {figures/quant/quant_swimmer_N8.dat};
\addplot[name path=swimmerN8b6hi, draw=none, forget plot] table[x=iteration, y=b6_hi] {figures/quant/quant_swimmer_N8.dat};
\addplot[colorC, opacity=0.15, forget plot] fill between[of=swimmerN8b6lo and swimmerN8b6hi];
\addplot[colorC, line width=0.8pt] table[x=iteration, y=b6_mean] {figures/quant/quant_swimmer_N8.dat};

\nextgroupplot[xmin=0, xmax=2000,
    ymin=0, ymax=3600,
    xtick distance=1000,
    ytick distance=1200,
    scaled y ticks=base 10:-3, ytick scale label code/.code={},
    ylabel={Hopper-v4\\return ($\times 10^3$)}]
\addplot[name path=hopperN2fp32lo, draw=none, forget plot] table[x=iteration, y=fp32_lo] {figures/quant/quant_hopper_N2.dat};
\addplot[name path=hopperN2fp32hi, draw=none, forget plot] table[x=iteration, y=fp32_hi] {figures/quant/quant_hopper_N2.dat};
\addplot[colorA, opacity=0.15, forget plot] fill between[of=hopperN2fp32lo and hopperN2fp32hi];
\addplot[colorA, line width=0.8pt] table[x=iteration, y=fp32_mean] {figures/quant/quant_hopper_N2.dat};
\addplot[name path=hopperN2b8lo, draw=none, forget plot] table[x=iteration, y=b8_lo] {figures/quant/quant_hopper_N2.dat};
\addplot[name path=hopperN2b8hi, draw=none, forget plot] table[x=iteration, y=b8_hi] {figures/quant/quant_hopper_N2.dat};
\addplot[colorB, opacity=0.15, forget plot] fill between[of=hopperN2b8lo and hopperN2b8hi];
\addplot[colorB, line width=0.8pt] table[x=iteration, y=b8_mean] {figures/quant/quant_hopper_N2.dat};
\addplot[name path=hopperN2b6lo, draw=none, forget plot] table[x=iteration, y=b6_lo] {figures/quant/quant_hopper_N2.dat};
\addplot[name path=hopperN2b6hi, draw=none, forget plot] table[x=iteration, y=b6_hi] {figures/quant/quant_hopper_N2.dat};
\addplot[colorC, opacity=0.15, forget plot] fill between[of=hopperN2b6lo and hopperN2b6hi];
\addplot[colorC, line width=0.8pt] table[x=iteration, y=b6_mean] {figures/quant/quant_hopper_N2.dat};

\nextgroupplot[xmin=0, xmax=2000,
    ymin=0, ymax=3600,
    xtick distance=1000,
    ytick distance=1200,
    scaled y ticks=base 10:-3, ytick scale label code/.code={}]
\addplot[name path=hopperN4fp32lo, draw=none, forget plot] table[x=iteration, y=fp32_lo] {figures/quant/quant_hopper_N4.dat};
\addplot[name path=hopperN4fp32hi, draw=none, forget plot] table[x=iteration, y=fp32_hi] {figures/quant/quant_hopper_N4.dat};
\addplot[colorA, opacity=0.15, forget plot] fill between[of=hopperN4fp32lo and hopperN4fp32hi];
\addplot[colorA, line width=0.8pt] table[x=iteration, y=fp32_mean] {figures/quant/quant_hopper_N4.dat};
\addplot[name path=hopperN4b8lo, draw=none, forget plot] table[x=iteration, y=b8_lo] {figures/quant/quant_hopper_N4.dat};
\addplot[name path=hopperN4b8hi, draw=none, forget plot] table[x=iteration, y=b8_hi] {figures/quant/quant_hopper_N4.dat};
\addplot[colorB, opacity=0.15, forget plot] fill between[of=hopperN4b8lo and hopperN4b8hi];
\addplot[colorB, line width=0.8pt] table[x=iteration, y=b8_mean] {figures/quant/quant_hopper_N4.dat};
\addplot[name path=hopperN4b6lo, draw=none, forget plot] table[x=iteration, y=b6_lo] {figures/quant/quant_hopper_N4.dat};
\addplot[name path=hopperN4b6hi, draw=none, forget plot] table[x=iteration, y=b6_hi] {figures/quant/quant_hopper_N4.dat};
\addplot[colorC, opacity=0.15, forget plot] fill between[of=hopperN4b6lo and hopperN4b6hi];
\addplot[colorC, line width=0.8pt] table[x=iteration, y=b6_mean] {figures/quant/quant_hopper_N4.dat};

\nextgroupplot[xmin=0, xmax=2000,
    ymin=0, ymax=3600,
    xtick distance=1000,
    ytick distance=1200,
    scaled y ticks=base 10:-3, ytick scale label code/.code={}]
\addplot[name path=hopperN8fp32lo, draw=none, forget plot] table[x=iteration, y=fp32_lo] {figures/quant/quant_hopper_N8.dat};
\addplot[name path=hopperN8fp32hi, draw=none, forget plot] table[x=iteration, y=fp32_hi] {figures/quant/quant_hopper_N8.dat};
\addplot[colorA, opacity=0.15, forget plot] fill between[of=hopperN8fp32lo and hopperN8fp32hi];
\addplot[colorA, line width=0.8pt] table[x=iteration, y=fp32_mean] {figures/quant/quant_hopper_N8.dat};
\addplot[name path=hopperN8b8lo, draw=none, forget plot] table[x=iteration, y=b8_lo] {figures/quant/quant_hopper_N8.dat};
\addplot[name path=hopperN8b8hi, draw=none, forget plot] table[x=iteration, y=b8_hi] {figures/quant/quant_hopper_N8.dat};
\addplot[colorB, opacity=0.15, forget plot] fill between[of=hopperN8b8lo and hopperN8b8hi];
\addplot[colorB, line width=0.8pt] table[x=iteration, y=b8_mean] {figures/quant/quant_hopper_N8.dat};
\addplot[name path=hopperN8b6lo, draw=none, forget plot] table[x=iteration, y=b6_lo] {figures/quant/quant_hopper_N8.dat};
\addplot[name path=hopperN8b6hi, draw=none, forget plot] table[x=iteration, y=b6_hi] {figures/quant/quant_hopper_N8.dat};
\addplot[colorC, opacity=0.15, forget plot] fill between[of=hopperN8b6lo and hopperN8b6hi];
\addplot[colorC, line width=0.8pt] table[x=iteration, y=b6_mean] {figures/quant/quant_hopper_N8.dat};

\nextgroupplot[xmin=0, xmax=3000,
    ymin=0, ymax=7500,
    xtick distance=1500,
    ytick distance=2500,
    scaled y ticks=base 10:-3, ytick scale label code/.code={},
    ylabel={Humanoid-v4\\return ($\times 10^3$)}]
\addplot[name path=humanoidN2fp32lo, draw=none, forget plot] table[x=iteration, y=fp32_lo] {figures/quant/quant_humanoid_N2.dat};
\addplot[name path=humanoidN2fp32hi, draw=none, forget plot] table[x=iteration, y=fp32_hi] {figures/quant/quant_humanoid_N2.dat};
\addplot[colorA, opacity=0.15, forget plot] fill between[of=humanoidN2fp32lo and humanoidN2fp32hi];
\addplot[colorA, line width=0.8pt] table[x=iteration, y=fp32_mean] {figures/quant/quant_humanoid_N2.dat};
\addplot[name path=humanoidN2b8lo, draw=none, forget plot] table[x=iteration, y=b8_lo] {figures/quant/quant_humanoid_N2.dat};
\addplot[name path=humanoidN2b8hi, draw=none, forget plot] table[x=iteration, y=b8_hi] {figures/quant/quant_humanoid_N2.dat};
\addplot[colorB, opacity=0.15, forget plot] fill between[of=humanoidN2b8lo and humanoidN2b8hi];
\addplot[colorB, line width=0.8pt] table[x=iteration, y=b8_mean] {figures/quant/quant_humanoid_N2.dat};
\addplot[name path=humanoidN2b6lo, draw=none, forget plot] table[x=iteration, y=b6_lo] {figures/quant/quant_humanoid_N2.dat};
\addplot[name path=humanoidN2b6hi, draw=none, forget plot] table[x=iteration, y=b6_hi] {figures/quant/quant_humanoid_N2.dat};
\addplot[colorC, opacity=0.15, forget plot] fill between[of=humanoidN2b6lo and humanoidN2b6hi];
\addplot[colorC, line width=0.8pt] table[x=iteration, y=b6_mean] {figures/quant/quant_humanoid_N2.dat};

\nextgroupplot[xmin=0, xmax=3000,
    ymin=0, ymax=7500,
    xtick distance=1500,
    ytick distance=2500,
    scaled y ticks=base 10:-3, ytick scale label code/.code={}]
\addplot[name path=humanoidN4fp32lo, draw=none, forget plot] table[x=iteration, y=fp32_lo] {figures/quant/quant_humanoid_N4.dat};
\addplot[name path=humanoidN4fp32hi, draw=none, forget plot] table[x=iteration, y=fp32_hi] {figures/quant/quant_humanoid_N4.dat};
\addplot[colorA, opacity=0.15, forget plot] fill between[of=humanoidN4fp32lo and humanoidN4fp32hi];
\addplot[colorA, line width=0.8pt] table[x=iteration, y=fp32_mean] {figures/quant/quant_humanoid_N4.dat};
\addplot[name path=humanoidN4b8lo, draw=none, forget plot] table[x=iteration, y=b8_lo] {figures/quant/quant_humanoid_N4.dat};
\addplot[name path=humanoidN4b8hi, draw=none, forget plot] table[x=iteration, y=b8_hi] {figures/quant/quant_humanoid_N4.dat};
\addplot[colorB, opacity=0.15, forget plot] fill between[of=humanoidN4b8lo and humanoidN4b8hi];
\addplot[colorB, line width=0.8pt] table[x=iteration, y=b8_mean] {figures/quant/quant_humanoid_N4.dat};
\addplot[name path=humanoidN4b6lo, draw=none, forget plot] table[x=iteration, y=b6_lo] {figures/quant/quant_humanoid_N4.dat};
\addplot[name path=humanoidN4b6hi, draw=none, forget plot] table[x=iteration, y=b6_hi] {figures/quant/quant_humanoid_N4.dat};
\addplot[colorC, opacity=0.15, forget plot] fill between[of=humanoidN4b6lo and humanoidN4b6hi];
\addplot[colorC, line width=0.8pt] table[x=iteration, y=b6_mean] {figures/quant/quant_humanoid_N4.dat};

\nextgroupplot[xmin=0, xmax=3000,
    ymin=0, ymax=7500,
    xtick distance=1500,
    ytick distance=2500,
    scaled y ticks=base 10:-3, ytick scale label code/.code={}]
\addplot[name path=humanoidN8fp32lo, draw=none, forget plot] table[x=iteration, y=fp32_lo] {figures/quant/quant_humanoid_N8.dat};
\addplot[name path=humanoidN8fp32hi, draw=none, forget plot] table[x=iteration, y=fp32_hi] {figures/quant/quant_humanoid_N8.dat};
\addplot[colorA, opacity=0.15, forget plot] fill between[of=humanoidN8fp32lo and humanoidN8fp32hi];
\addplot[colorA, line width=0.8pt] table[x=iteration, y=fp32_mean] {figures/quant/quant_humanoid_N8.dat};
\addplot[name path=humanoidN8b8lo, draw=none, forget plot] table[x=iteration, y=b8_lo] {figures/quant/quant_humanoid_N8.dat};
\addplot[name path=humanoidN8b8hi, draw=none, forget plot] table[x=iteration, y=b8_hi] {figures/quant/quant_humanoid_N8.dat};
\addplot[colorB, opacity=0.15, forget plot] fill between[of=humanoidN8b8lo and humanoidN8b8hi];
\addplot[colorB, line width=0.8pt] table[x=iteration, y=b8_mean] {figures/quant/quant_humanoid_N8.dat};
\addplot[name path=humanoidN8b6lo, draw=none, forget plot] table[x=iteration, y=b6_lo] {figures/quant/quant_humanoid_N8.dat};
\addplot[name path=humanoidN8b6hi, draw=none, forget plot] table[x=iteration, y=b6_hi] {figures/quant/quant_humanoid_N8.dat};
\addplot[colorC, opacity=0.15, forget plot] fill between[of=humanoidN8b6lo and humanoidN8b6hi];
\addplot[colorC, line width=0.8pt] table[x=iteration, y=b6_mean] {figures/quant/quant_humanoid_N8.dat};
\end{groupplot}
\end{tikzpicture}

  \caption{Return over training, unquantized (\texttt{float32}) and under
  stochastic quantization at 8 and 6 bits. Rows: environments; columns: number of agents
  $N$. Lines are the mean over five seeds, smoothed with a centred 21-iteration
  rolling mean; bands span the minimum to the maximum across seeds. Hopper
  uses its own configuration (Section~\ref{sec:implementation}) in all three
  series, Swimmer and Humanoid the original one.}
  \label{fig:quant}
\end{figure}

\begin{figure}[tbp]
  \centering
  %% Uplink of one whole training run at N=8, summed over agents, by bit width.
%% Measured cumulative_comm_bytes, identical across seeds. Log scale. The
%% unquantized run is shown as metered at the paper's 64 bits per value and
%% at the float32 it actually exchanges (half of it).
%% Data: figures/quant/quant_overhead.dat (scripts/export_quant_curves.py).
\begin{tikzpicture}
\begin{axis}[
    ybar, bar width=8pt,
    width=0.62\textwidth, height=0.38\textwidth,
    symbolic x coords={Swimmer,Hopper,Humanoid},
    xtick=data,
    enlarge x limits=0.25,
    ymode=log, ymin=0.01, ymax=5000,
    log origin=infty,
    ylabel={Uplink per run at $N=8$ (GB)},
    ymajorgrids=true, grid style=dashed,
    tick label style={font=\footnotesize},
    label style={font=\footnotesize},
    point meta=rawy,
    nodes near coords={\pgfmathprintnumber[fixed, precision=2]{\pgfplotspointmeta}},
    nodes near coords style={font=\tiny, rotate=90, anchor=west},
    legend style={font=\footnotesize, legend columns=2, draw=none, cells={anchor=west},
                  at={(0.5,1.02)}, anchor=south,
                  /tikz/every even column/.append style={column sep=0.8em}},
    legend image code/.code={\draw[#1] (0cm,-0.1cm) rectangle (0.2cm,0.12cm);},
]
\addplot[draw=colorA, dashed, fill=colorA!15] table[x=env, y=b64] {figures/quant/quant_overhead.dat};
\addlegendentry{Unquantized, 64-bit accounting (original paper)}
\addplot[draw=colorA, fill=colorA!60] table[x=env, y=fp32] {figures/quant/quant_overhead.dat};
\addlegendentry{Unquantized, float32}
\addplot[draw=colorB, fill=colorB!60] table[x=env, y=b8] {figures/quant/quant_overhead.dat};
\addlegendentry{8 bits}
\addplot[draw=colorC, fill=colorC!60] table[x=env, y=b6] {figures/quant/quant_overhead.dat};
\addlegendentry{6 bits}
\end{axis}
\end{tikzpicture}

  \caption{Uplink of one complete training run at $N=8$, summed over agents;
  log scale. The unquantized run is shown both as metered, at the original paper's 64
  bits per value, and at the \texttt{float32} it actually exchanges, half of
  that. Measured byte counts, identical across seeds. Eight and six bits cut
  the uplink by $4\times$ and $5.3\times$ against \texttt{float32}, or $8\times$
  and $10.7\times$ against the 64-bit accounting, on every environment; the
  totals scale linearly with $N$.}
  \label{fig:quantoverhead}
\end{figure}

\paragraph{Eight bits preserve the result.} On Swimmer, 8-bit and unquantized
runs are indistinguishable at every agent count. On Humanoid, 8 bits stays
within one standard deviation of the unquantized result at $N=2$ and $N=4$
and finishes $5\%$ lower at $N=8$ ($5997 \pm 259$ against $6329 \pm 334$),
also about one standard deviation. Learning still speeds up with the number of
agents, which is the paper's claim: at iteration 1000 the
smoothed Humanoid return at 8 bits is $580$, $905$ and $3428$ for $N = 2, 4,
8$, against $602$, $944$ and $3863$ unquantized, and the same ordering holds at
every checkpoint on both environments. Eight bits therefore removes a further
factor of four from the uplink against the \texttt{float32} values actually
exchanged (eight at the paper's 64-bit accounting), on top of the
$\Oh(d^2) \to \Oh(d)$ reduction, without giving up either the final reward or the scaling with $N$.

\paragraph{Six bits are too few.} On Humanoid, 6 bits collapses at every agent
count, to $497$, $676$ and $1514$. Five seeds put this beyond seed variance and
confirm the threshold the single-seed scan found (5543 at 7 bits, 1991 at
6). The ordering in $N$ survives, at about a quarter of the
return. Swimmer's final values hold at 6 bits, but the cost shows elsewhere:
the seed spread at $N=2$ and $N=4$ grows from about $1$ to $7.4$ and $3.4$,
and early in training four agents learn \emph{slower} than two (smoothed return
$75$ against $95$ at iteration 100), so the faster-with-more-agents ordering
is lost before it is recovered at convergence.

\paragraph{Where the threshold sits.} The single-seed scan at $N=8$ shows
a sharp threshold rather than a gradual decay. Swimmer is flat from 6 to 32
bits and collapses to $94.4$ at 5; Humanoid holds $5542.6$ at 7 bits and falls
to $1990.5$ at 6. The threshold moves in the direction $d$ does, which is what
an unbiased quantizer predicts. The variance injected per element is bounded by
$s^2/4$ and enters a $d$-dimensional sum, so Humanoid's policy
($d \approx 7.3\times10^{5}$) needs one bit more than Swimmer's
($d = \num{4868}$) before the aggregated direction stops degrading. Two
environments do not test that reading, and a scan on a fourth would. Taken
together with the five-seed grid, 8 bits is the operating point that holds on
every environment we tested; 7 bits on Humanoid and 6 on Swimmer are
plausible but rest on single seeds or show degraded early learning.

\paragraph{Hopper.} Hopper tolerates quantization better than Humanoid. At
$N=8$ both bit widths finish within the seed spread of the unquantized runs
($2581 \pm 466$ at 8 bits and $2590 \pm 338$ at 6, against $2656 \pm 286$),
and 6 bits does not collapse as it does on Humanoid. At $N=4$ both quantized
runs finish lower ($2197$ and $2096$ against $2536$), by about one standard
deviation at 8 bits and more at 6; at $N=2$ they finish higher. The final
reward still increases with $N$ at both bit widths ($2052$, $2197$, $2581$ at
8 bits; $1941$, $2096$, $2590$ at 6), so the scaling with the number of agents
survives quantization here too.

\subsection{What the byte counts omit}

The reported bytes are $d \cdot b / 8$ per transmission and exclude the
quantization metadata a real deployment would also send: $\min x$ and the scale
for the stochastic scheme (two \texttt{float32}, 64 bits per transmission).
Relative to a $d \cdot b$-bit payload the
overhead is $64/(d\,b)$, so roughly $1$--$3\%$ in the worst case here ($b=1$, $d$
of a few thousand) and below $0.5\%$ at 8 bits. The downlink broadcast of
$\bthe$ and $\by$ is not counted either, for the same reason it is not counted
in the paper. Neither omission changes the ordering between bit widths, but
both make the reported compression slightly optimistic.

\paragraph{Quantizing the baseline instead.} A symmetric comparison would
quantize standard-average FedNPG's $\bH_i$ as well. It does not change the
conclusion, and the arithmetic shows why without an experiment. At $b$ bits the
standard method uplinks about $d^2 b$ bits per agent per round, against
$2 \cdot 64\, d$ for unquantized FedNPG-ADMM, so the two break even only at
$b \approx 128/d$, or $64/d$ against \texttt{float32}. On Swimmer that is at
most $0.03$ bits, and smaller still on the
larger policies: no bit width makes the $\Oh(d^2)$ method competitive, and
quantizing FedNPG-ADMM widens the gap further.
